\section{Length}
\label{section-length}

\begin{definition}
\label{definition-length}
Let $R$ be a ring. For any $R$-module $M$
we define the {\it length} of $M$ over $R$ by the
formula
$$
\text{length}_R(M)
=
\sup
\{
n \in \mathbf Z_{\geq0}
\mid
\exists\ 0 = M_0 \subset M_1 \subset \ldots \subset M_n = M,
\text{ }M_i \not = M_{i + 1}\quad(0\leq i<n)
\}.
$$
\end{definition}

\noindent
The supremum is taken in $\mathbf Z_{\geq0}\cup\{\infty\}$.
For $M=0$ the length-zero chain is allowed, so its length is zero.
In other words it is the supremum of the lengths of chains
of submodules. There is an obvious notion of when a chain
of submodules is a refinement of another. This gives a
partial ordering on the collection of all chains of submodules,
with the smallest chain having the shape $0 = M_0 \subset M_1 = M$
if $M$ is not zero.
We note the obvious fact that if the length of
$M$ is finite, then every chain can be refined to a
maximal chain. But it is not as obvious that all maximal
chains have the same length (as we will see later).

\begin{lemma}
\label{lemma-finite-length-finite}
\begin{slogan}
Modules of finite length are finite.
\end{slogan}
Let $R$ be a ring.
Let $M$ be an $R$-module.
If $\text{length}_R(M) < \infty$ then $M$ is a finite $R$-module.
\end{lemma}

\begin{proof}
Write $\ell=\operatorname{length}_R(M)<\infty$.
If $M=0$, its empty generating list suffices. Otherwise start with
$N_0=0$. Whenever $N_j\neq M$, choose
$x_{j+1}\in M\setminus N_j$ and put
$N_{j+1}=N_j+Rx_{j+1}$. Each such step is strict.
There cannot be $\ell+1$ such steps: adjoining the final endpoint
$M$ if it is not already present would give a chain from zero
to $M$ with at least $\ell+1$ strict steps, contrary to the
definition of length. Therefore $N_j=M$ for some $j\leq\ell$,
and $x_1,\ldots,x_j$ generate the original $M$.

The same bound applies to every subquotient $N/P$ of $M$.
Indeed a strict chain in $N/P$ pulls back under the original
quotient map $N\to N/P$ to a strict chain from $P$ to $N$.
Adjoin zero if $P\neq0$ and $M$ if $N\neq M$. This is a
chain of submodules of $M$, so the original chain has at most
$\ell$ strict steps. Thus $N/P$ has length at most $\ell$
and is generated by at most its own length many elements by the
same construction. Moreover, no ascending or descending chain
of submodules of $M$ can have more than $\ell$ strict changes:
any finite selection of its strict changes, ordered increasingly
and completed with the endpoints, gives such a chain in $M$.
In particular both chain conditions hold for $M$.
\end{proof}

\begin{lemma}
\label{lemma-length-additive}
\begin{slogan}
Length is additive in short exact sequences.
\end{slogan}
If $0 \to M' \to M \to M'' \to 0$
is a short exact sequence of modules over $R$ then
the length of $M$ is the sum of the
lengths of $M'$ and $M''$.
\end{lemma}

\begin{proof}
Write the original sequence as
$$
0\longrightarrow M'\xrightarrow{u}M\xrightarrow{q}M''
\longrightarrow0.
$$
Lengths take values in $\mathbf Z_{\geq0}\cup\{\infty\}$,
with $a+\infty=\infty+a=\infty$ for every such $a$.
Given strict chains from zero to $M'$ and $M''$, of lengths
$a$ and $b$, respectively, first apply $u$ to the chain in
$M'$. Its last member is $u(M')=\ker(q)$.
Continue with the inverse images under $q$ of the nonzero steps
of the chain in $M''$. Injectivity of $u$ and surjectivity of
$q$ show that every indicated step is strict. This yields a
chain in $M$ of length $a+b$. If either length on the right is
infinite, finite chains of arbitrarily large length give
$\operatorname{length}_R(M)=\infty$. If both are finite, each
is an attained integer supremum, so choosing chains with those
lengths proves
$\operatorname{length}_R(M)\geq
\operatorname{length}_R(M')+\operatorname{length}_R(M'')$.
This construction includes a zero module with its length-zero chain.

Conversely, let $0=M_0\subsetneq\cdots\subsetneq M_n=M$
be a finite strict chain. Put
$M_i'=u^{-1}(M_i)$ and $M_i''=q(M_i)$, using the original
maps rather than silently identifying $M'$ with a different module.
Let $a$ and $b$ count only the strict changes in these two
nondecreasing sequences. Removing their repeated terms gives
chains of lengths $a$ and $b$ in $M'$ and $M''$.
At least one changes at each of the $n$ original steps.
For if $M_i'=M_{i+1}'$ and $M_i''=M_{i+1}''$, take
$x\in M_{i+1}$. There is $y\in M_i$ with $q(x)=q(y)$.
Then $x-y=u(z)$ for some $z\in M'$ and
$u(z)\in M_{i+1}$, so $z\in M_{i+1}'=M_i'$.
Consequently $u(z)\in M_i$ and $x=y+u(z)\in M_i$,
contradicting strictness. It follows that
$$
n\leq a+b\leq
\operatorname{length}_R(M')+\operatorname{length}_R(M'').
$$
Taking the supremum over the original chains proves the reverse
inequality, including infinite lengths. This proves additivity.
\end{proof}

\begin{lemma}
\label{lemma-length-infinite}
Let $R$ be a local ring with maximal ideal $\mathfrak m$.
If $M$ is an $R$-module and $\mathfrak m^n M \not = 0$ for all $n \geq 0$,
then $\text{length}_R(M) = \infty$. In other words, if $M$
has finite length then $\mathfrak m^nM = 0$ for some $n$.
\end{lemma}

\begin{proof}
Fix an integer $n\geq0$ such that $\mathfrak m^nM\neq0$.
A nonzero member of this module is a finite sum of products
$f_1\cdots f_nx$, with every $f_i\in\mathfrak m$ and $x\in M$.
At least one of these products is nonzero; choose it and put
$y_j=f_1\cdots f_jx$ for $0\leq j\leq n$, where $y_0=x$.
All $y_j$ are nonzero, since multiplying a zero earlier product
by the remaining original factors would give $y_n=0$.
The inclusions
$$
0\subsetneq Ry_n\subsetneq Ry_{n-1}\subsetneq\cdots
\subsetneq Ry_0\subset M
$$
are strict up to $Ry_0$. To verify any step with $1\leq j\leq n$,
suppose $Ry_{j-1}=Ry_j$. Then
$y_{j-1}=g y_j=g f_j y_{j-1}$ for some $g\in R$, hence
$(1-gf_j)y_{j-1}=0$. Since $gf_j\in\mathfrak m$, the original
factor $1-gf_j$ is a unit in the local ring. This forces
$y_{j-1}=0$, a contradiction. The first inclusion is strict
because $y_n\neq0$. Append $M$ if necessary; if $Ry_0=M$,
there is no extra final step. Thus
$\operatorname{length}_R(M)\geq n+1$.
For $n=0$ this uses the nonzero element $x$ and the empty product,
and gives the same bound.

If every $\mathfrak m^nM$ is nonzero, these bounds for all $n$
give infinite length. They also give the sharper quantitative
conclusion
$$
\mathfrak m^\ell M=0
\quad\text{if}\quad
\ell=\operatorname{length}_R(M)<\infty.
$$
Indeed a nonzero left-hand module would imply
$\ell\geq\ell+1$. When $\ell=0$, the definition of length
already gives $M=0$, so the displayed assertion remains valid.
\end{proof}

\begin{lemma}
\label{lemma-length-independent}
Let $R \to S$ be a ring map. Let $M$ be an $S$-module.
We always have $\text{length}_R(M) \geq \text{length}_S(M)$.
If $R \to S$ is surjective then equality holds.
\end{lemma}

\begin{proof}
Every $S$-submodule of the original $M$ is an $R$-submodule by
restriction of scalars. Hence every strict chain for the former
action gives the identical strict chain for the latter, proving
the inequality even for infinite length.
If $R\to S$ is surjective, then for any $R$-submodule $N\subset M$,
$s\in S$ and $x\in N$, choose $r\in R$ mapping to $s$.
The equality $sx=rx$ puts $sx$ in $N$. Thus the two collections
of submodules, and all their chains, agree exactly.

The proof gives a weaker sufficient hypothesis for equality:
surjectivity of the original composite
$R\to S/\operatorname{Ann}_S(M)$ suffices.
For then any $s\in S$ has an $r\in R$ such that their images
in this quotient agree; their difference annihilates every
original $x\in M$, so once again $sx=rx$.
The original actions are retained and their equality on each
element proves equality of the submodule lattices. When $M=0$,
both lattices have a single member and both lengths are zero.
\end{proof}

\begin{lemma}
\label{lemma-dimension-is-length}
Let $R$ be a ring with maximal ideal $\mathfrak m$.
Suppose that $M$ is an $R$-module with
$\mathfrak m M  =  0$. Then the length of $M$ as
an $R$-module agrees with its dimension as a vector space over
$R/\mathfrak m$, interpreting every infinite dimension as $\infty$.
The length is finite if and only if $M$ is a finite $R$-module.
\end{lemma}

\begin{proof}
Set $k=R/\mathfrak m$. The original action on $M$ factors
through $k$, since $\mathfrak mM=0$; explicitly
$(r+\mathfrak m)x=rx$, which is independent of the representative.
An $R$-submodule is exactly a $k$-vector subspace by
Lemma \ref{lemma-length-independent}.
If $\dim_k M=d<\infty$, a strict inclusion of vector subspaces
increases dimension by at least one. Thus every strict chain has
at most $d$ steps. Given a basis $x_1,\ldots,x_d$, the spans
of its first $i$ vectors, for $0\leq i\leq d$, give a chain
of exactly $d$ steps. Therefore the length is $d$, including $d=0$.
If the dimension is infinite, for each finite $n$ choose
$n$ independent vectors. Their initial spans form $n$ strict
steps, and adjoining $M$ as endpoint if needed gives length at
least $n$. The length is therefore $\infty$; it is this extended
integer value, not the particular infinite basis cardinal, that
is asserted here.
Finally, a finite $R$-generating list is exactly a finite
$k$-spanning list because the two actions just described have the
same spans. Such a list contains a basis after removing dependent
vectors, and a finite basis is itself a finite $R$-generating
list. This proves the final equivalence.
\end{proof}

\begin{lemma}
\label{lemma-length-localize}
Let $R$ be a ring. Let $M$ be an $R$-module. Let $S \subset R$ be
a multiplicative subset. Then
$\text{length}_R(M) \geq \text{length}_{S^{-1}R}(S^{-1}M)$.
\end{lemma}

\begin{proof}
Let $j:M\to S^{-1}M$ be the original map $x\mapsto x/1$.
For an $S^{-1}R$-submodule $L\subset S^{-1}M$, put
$j^{-1}(L)=\{x\in M:x/1\in L\}$. Its localization has image
exactly $L$. Indeed if $x/s\in L$, then $x/1=s(x/s)\in L$,
so $x\in j^{-1}(L)$; conversely $x/1\in L$ implies
$x/s=(1/s)(x/1)\in L$. Localization of the inclusion
$j^{-1}(L)\subset M$ is injective, so this identifies the actual
localized submodule with $L$.
Consequently a strict chain $0=L_0\subsetneq\cdots\subsetneq
L_n=S^{-1}M$ pulls back to a strict chain
$j^{-1}(L_0)\subsetneq\cdots\subsetneq j^{-1}(L_n)=M$:
equality of two adjacent inverse images would give equality of
their localized images. The first inverse image is $\ker(j)$,
which need not be zero. Adjoin zero when it is nonzero. This
produces at least $n$ strict steps from zero to $M$, proving
the length inequality. If $S^{-1}M=0$, its length is zero and
the inequality holds directly; no nonzero localization is assumed.
For finite length, the proof of Lemma \ref{lemma-simple-pieces}
below computes the exact loss: it is the number of original
composition factors whose maximal ideals meet $S$, counted
with their multiplicities.
\end{proof}

\begin{lemma}
\label{lemma-length-finite}
Let $R$ be a ring with finitely generated
maximal ideal $\mathfrak m$. (For example, this holds if $R$ is Noetherian.)
Suppose that $M$ is a finite $R$-module with
$\mathfrak m^n M  =  0$ for some integer $n\geq0$.
Then $\text{length}_R(M) < \infty$.
\end{lemma}

\begin{proof}
If $n=0$, then $M=\mathfrak m^0M=0$ and the claim follows.
Assume $n\geq1$ and write $k=R/\mathfrak m$.
Choose original generators $x_1,\ldots,x_r$ of $M$ and
$f_1,\ldots,f_t$ of $\mathfrak m$. For $i\geq0$, put
$$
E_i=\{(e_1,\ldots,e_t)\in\mathbf Z_{\geq0}^{,t}:
e_1+\cdots+e_t=i\}.
$$
The quotient $\mathfrak m^iM/\mathfrak m^{i+1}M$ is killed
by $\mathfrak m$ and has the following explicit surjection
of $k$-vector spaces:
$$
\begin{aligned}
k^{\{1,\ldots,r\}\times E_i}&\longrightarrow
\mathfrak m^iM/\mathfrak m^{i+1}M,\\
(\bar a_{j,e})_{j,e}&\longmapsto
\sum_{j=1}^r\sum_{e\in E_i}
a_{j,e} f_1^{e_1}\cdots f_t^{e_t}x_j
\pmod{\mathfrak m^{i+1}M}.
\end{aligned}
$$
Changing any representative $a_{j,e}$ by an element of
$\mathfrak m$ changes the indicated summand by a member of
$\mathfrak m^{i+1}M$, so the map is well defined. The products
in the specified ideal generators generate $\mathfrak m^i$;
expanding a member of $\mathfrak m^iM$ in these products and
the original $x_j$ proves surjectivity. Thus each layer has
finite dimension over $k$, and its length is that dimension by
Lemma \ref{lemma-dimension-is-length}. Applying additivity
successively to the original filtration by powers gives the exact formula
$$
\operatorname{length}_R(M)
=\sum_{i=0}^{n-1}\dim_k
 (\mathfrak m^iM/\mathfrak m^{i+1}M)
\leq r\sum_{i=0}^{n-1}\#E_i<\infty.
$$
Repeated adjacent powers contribute their actual zero-dimensional
quotients. If $t\geq1$, distributing $i$ positions among $t$
nonnegative exponents gives $\#E_i=\binom{t+i-1}{i}$: encode
the tuple by $i$ identical positions separated into $t$ blocks
by $t-1$ separators, choosing the $i$ positions among the
$t+i-1$ total positions. If $t=0$, there is one empty tuple
in $E_0$ and no tuple in $E_i$ for $i>0$.

In fact finite generation of $\mathfrak m$ can be replaced by
finite dimension of the original cotangent space
$\mathfrak m/\mathfrak m^2$ over $k$.
Choose lifts $f_1,\ldots,f_t\in\mathfrak m$ of a finite
spanning list and put $J=(f_1,\ldots,f_t)$.
Then $\mathfrak m=J+\mathfrak m^2$.
For $i\geq1$, expanding this actual ideal equality retains all terms:
$$
\mathfrak m^i
=(J+\mathfrak m^2)^i
=\sum_{h=0}^i J^{i-h}\mathfrak m^{2h}.
$$
Every term with $h\geq1$ lies in $\mathfrak m^{i+h}$,
hence in $\mathfrak m^{i+1}$. Conversely both $J^i$ and
$\mathfrak m^{i+1}$ lie in $\mathfrak m^i$. Therefore
$\mathfrak m^i=J^i+\mathfrak m^{i+1}$ and
$\mathfrak m^iM=J^iM+\mathfrak m^{i+1}M$.
All the terms with $h\geq1$ have zero image under this particular
quotient map; no equality $\mathfrak m=J$ has been assumed.
The same displayed coefficient map is consequently surjective
for every layer, with these original lifts $f_j$.
It proves the same finite length conclusion and the same explicit
bound whenever $M$ is finite and $\mathfrak m^nM=0$.
This argument covers $t=0$ as well: then
$\mathfrak m^iM=\mathfrak m^{i+1}M$ for $i\geq1$, so
nilpotence on $M$ forces all those layers and $\mathfrak mM$ to vanish.
\end{proof}

\begin{definition}
\label{definition-simple-module}
Let $R$ be a ring. Let $M$ be an $R$-module.
We say $M$ is {\it simple} if $M \not = 0$ and
every submodule of $M$ is either equal to $M$ or
to $0$.
\end{definition}

\begin{lemma}
\label{lemma-characterize-length-1}
Let $R$ be a ring. Let $M$ be an $R$-module.
The following are equivalent:
\begin{enumerate}
\item $M$ is simple,
\item $\text{length}_R(M) = 1$, and
\item $M \cong R/\mathfrak m$ for some maximal ideal
$\mathfrak m \subset R$.
\end{enumerate}
\end{lemma}

\begin{proof}
Let $\mathfrak m$ be a maximal ideal of $R$.
By Lemma \ref{lemma-dimension-is-length} the module
$R/\mathfrak m$ has length $1$. The equivalence of
the first two assertions is tautological.
Suppose that $M$ is simple. Choose $x \in M$, $x \not = 0$.
As $M$ is simple we have $M = R \cdot x$.
Let $I \subset R$ be the annihilator of $x$, i.e.,
$I = \{f \in R \mid fx = 0\}$. The map $R/I \to M$,
$f \bmod I \mapsto fx$ is an isomorphism, hence
$R/I$ is a simple $R$-module. Since $R/I \not = 0$ we see $I \not = R$.
Let $\mathfrak m$ be a maximal ideal containing $I$.
If $I \not = \mathfrak m$, then $\mathfrak m /I \subset R/I$
is a nontrivial submodule contradicting the simplicity
of $R/I$. Hence we see $I = \mathfrak m$ as desired.
\end{proof}

\begin{lemma}
\label{lemma-simple-pieces}
Let $R$ be a ring. Let $M$ be a finite length $R$-module.
Choose any maximal chain of submodules
$$
0 = M_0 \subset M_1 \subset M_2 \subset \ldots \subset M_n = M
$$
with $M_i \not = M_{i-1}$, $i = 1, \ldots, n$. Then
\begin{enumerate}
\item $n = \text{length}_R(M)$,
\item each $M_i/M_{i-1}$ is simple,
\item each $M_i/M_{i-1}$ is of the form
$R/\mathfrak m_i$ for some maximal ideal $\mathfrak m_i$,
\item given a maximal ideal $\mathfrak m \subset R$
we have
$$
\# \{i \mid \mathfrak m_i = \mathfrak m\}
=
\text{length}_{R_{\mathfrak m}} (M_{\mathfrak m}).
$$
\end{enumerate}
\end{lemma}

\begin{proof}
For each $i$, let $q_i:M_i\to M_i/M_{i-1}$ be the original
quotient map. If its nonzero target were not simple, a nonzero
proper submodule would have inverse image strictly between
$M_{i-1}$ and $M_i$. Inserting it would contradict maximality
of the given chain. Thus every quotient is simple, and
Lemma \ref{lemma-characterize-length-1} gives an isomorphism
$M_i/M_{i-1}\cong R/\mathfrak m_i$ and length one.
Applying Lemma \ref{lemma-length-additive} to the successive
short exact sequences gives $\operatorname{length}_R(M)=n$.
For $M=0$ the chain has no strict steps, the lists of quotients
and maximal ideals are empty, and this equality reads $0=0$.

For a maximal ideal $\mathfrak m$, localization gives the exact sequences
$$
0\longrightarrow(M_{i-1})_{\mathfrak m}
\longrightarrow(M_i)_{\mathfrak m}
\longrightarrow(M_i/M_{i-1})_{\mathfrak m}
\longrightarrow0.
$$
If $\mathfrak m_i\neq\mathfrak m$, choose
$a\in\mathfrak m_i\setminus\mathfrak m$; such an element exists
because distinct maximal ideals cannot contain one another.
It annihilates $R/\mathfrak m_i$ and is inverted at
$\mathfrak m$, so this quotient localizes to zero.
If $\mathfrak m_i=\mathfrak m$, the isomorphism with the
residue field of the original local ring is explicitly
$$
(R/\mathfrak m)_{\mathfrak m}
\longrightarrow R_{\mathfrak m}/\mathfrak mR_{\mathfrak m},
\qquad
(r+\mathfrak m)/s\longmapsto r/s+\mathfrak mR_{\mathfrak m}.
$$
Its inverse sends $r/s+\mathfrak mR_{\mathfrak m}$ to
$(r+\mathfrak m)/s$. The localization fraction relations hold
on both sides, and coefficients in $\mathfrak m$ act as zero,
so both maps are well defined; their composites are identities.
The resulting nonzero quotient has length one over $R_{\mathfrak m}$.
Additivity proves exactly the asserted count.

More generally let $S\subset R$ be any multiplicative subset.
For each factor $R/\mathfrak m_i$, its localization is zero
if $S\cap\mathfrak m_i\neq\varnothing$. Otherwise every
$s\in S$ has nonzero, hence invertible, image in the field
$R/\mathfrak m_i$, and
$$
S^{-1}(R/\mathfrak m_i)\longrightarrow R/\mathfrak m_i,
\qquad (r+\mathfrak m_i)/s\longmapsto
(r+\mathfrak m_i)(s+\mathfrak m_i)^{-1}
$$
is an isomorphism, with inverse $z\mapsto z/1$.
It identifies this localized module with the simple quotient
$S^{-1}R/S^{-1}\mathfrak m_i$.
Localizing every short exact sequence and adding the lengths gives
$$
\operatorname{length}_{S^{-1}R}(S^{-1}M)
=\#\{i: S\cap\mathfrak m_i=\varnothing\},
$$
and therefore the exact loss of length is
$\#\{i:S\cap\mathfrak m_i\neq\varnothing\}$.
This refines Lemma \ref{lemma-length-localize} for finite length.
For an arbitrary prime $\mathfrak p$, the same calculation with
$S=R\setminus\mathfrak p$ retains a factor precisely when
$\mathfrak m_i\subset\mathfrak p$, hence precisely when
$\mathfrak p=\mathfrak m_i$.
Thus the original support of $M$ is the finite set of distinct
$\mathfrak m_i$, and
$$
\operatorname{length}_R(M)
=\sum_{\mathfrak m\in\operatorname{Supp}_R(M)}
  \operatorname{length}_{R_{\mathfrak m}}(M_{\mathfrak m}).
$$
In particular the multiplicity of each simple module
$R/\mathfrak m$ is independent of the chosen maximal chain.
The ideal $\mathfrak m$ is determined by that module, since
its annihilator is exactly $\mathfrak m$; hence these counts
determine the multiset of simple factors, not only the number
of steps.

The two chain conditions also characterize finite length for
an arbitrary $R$-module. Necessity was proved in the proof of
Lemma \ref{lemma-finite-length-finite}.
Conversely assume both conditions. In every nonzero submodule
$L$, the ascending chain condition supplies a maximal proper
submodule: if none existed, starting from zero one could choose
an infinite strictly increasing sequence of proper submodules
of $L$. Successively choose such a maximal proper submodule
inside each nonzero module, starting from $M$.
The descending chain condition forces this sequence to reach
zero after finitely many strict steps. Reversing it gives a
finite chain whose quotients are simple. Each quotient has
length one by Lemma \ref{lemma-characterize-length-1}, so
additivity gives finite length for the original $M$.
\end{proof}

\begin{lemma}
\label{lemma-pushdown-module}
Let $A$ be a local ring with maximal ideal $\mathfrak m$.
Let $B$ be a semi-local ring with maximal ideals $\mathfrak m_i$,
$i = 1, \ldots, n$.
Suppose that $A \to B$ is a homomorphism such that each $\mathfrak m_i$
lies over $\mathfrak m$ and such that
$$
[\kappa(\mathfrak m_i) : \kappa(\mathfrak m)] < \infty.
$$
Let $M$ be a $B$-module of finite length.
Then
$$
\text{length}_A(M) = \sum\nolimits_{i = 1, \ldots, n}
[\kappa(\mathfrak m_i) : \kappa(\mathfrak m)]
\text{length}_{B_{\mathfrak m_i}}(M_{\mathfrak m_i}),
$$
in particular $\text{length}_A(M) < \infty$.
\end{lemma}

\begin{proof}
Choose a maximal chain of $B$-submodules of the original $M$
as in Lemma \ref{lemma-simple-pieces}. Each of its quotients is
isomorphic to $B/\mathfrak m_{i(j)}=\kappa(\mathfrak m_{i(j)})$.
The original map $A\to B$ induces an injective field map
$\kappa(\mathfrak m)=A/\mathfrak m\to B/\mathfrak m_i$:
its kernel before quotienting is exactly the contraction
$\mathfrak m_i\cap A=\mathfrak m$.
The $A$-action on this quotient is the action through that map.
By Lemma \ref{lemma-dimension-is-length}, its $A$-length is
exactly $[\kappa(\mathfrak m_i):\kappa(\mathfrak m)]$.
Additivity along the original chain, now regarded as a chain
of $A$-submodules, gives the sum of these degrees over its factors.
Lemma \ref{lemma-simple-pieces} identifies the number of factors
at $\mathfrak m_i$ with
$\operatorname{length}_{B_{\mathfrak m_i}}(M_{\mathfrak m_i})$.
Grouping the finite sum by those maximal ideals proves the
displayed formula and its finiteness. When $M=0$, the chain
is empty and all terms are zero, as required.

The assumptions need only concern the maximal ideals in the
support of $M$. More precisely, let $\alpha:A\to B$ be any
ring homomorphism and let $M$ have finite length over $B$.
The preceding composition-factor proof shows that
$Q=\operatorname{Supp}_B(M)$ is a finite set of maximal ideals.
For each $\mathfrak q\in Q$, put
$\mathfrak p_{\mathfrak q}=\alpha^{-1}(\mathfrak q)$.
Then $M$ has finite length over $A$ if and only if every
$\mathfrak p_{\mathfrak q}$ is maximal in $A$ and every
extension $\kappa(\mathfrak q)/\kappa(\mathfrak p_{\mathfrak q})$
has finite degree. Under these conditions the exact formula is
$$
\operatorname{length}_A(M)
=\sum_{\mathfrak q\in Q}
[\kappa(\mathfrak q):\kappa(\mathfrak p_{\mathfrak q})]
\operatorname{length}_{B_{\mathfrak q}}(M_{\mathfrak q}).
$$
To prove the criterion in both directions, fix one composition
factor $k=B/\mathfrak q$. The original map identifies
$D=A/\mathfrak p_{\mathfrak q}$ with a subring of the field $k$.
If $\mathfrak p_{\mathfrak q}$ is maximal, $D$ is its residue
field and the $A$-length of $k$ is its dimension over $D$ by
Lemma \ref{lemma-dimension-is-length}, finite exactly in the
stated finite-degree case.
If $\mathfrak p_{\mathfrak q}$ is not maximal, the domain $D$
is not a field. Choose a nonzero nonunit $a\in D$.
Its image is nonzero in $k$, and the original $A$-submodules
$$
D\subsetneq a^{-1}D\subsetneq a^{-2}D\subsetneq\cdots\subset k
$$
are strictly increasing. Indeed equality of two consecutive
members would give $a^{-(j+1)}=a^{-j}d$ for some $d\in D$,
and multiplying in $k$ by $a^{j+1}$ would yield $1=ad$ in
the embedded $D$, contradicting the choice of $a$.
Adjoining zero and the endpoint $k$ gives finite chains of
arbitrarily large length, so $\operatorname{length}_A(k)=\infty$.
Finally, the original finite composition chain in $M$ remains
an exact filtration of $A$-modules. Additivity shows that its
$A$-length is the sum of its factor lengths, and is finite
exactly when each occurring factor has finite $A$-length.
Every $\mathfrak q\in Q$ occurs with the strictly positive
multiplicity specified above. This proves necessity, sufficiency
and the formula without any local or semilocal hypothesis on
the whole rings. If $Q$ is empty, $M=0$ and the criterion is
vacuous with an empty sum equal to zero.
\end{proof}

\begin{lemma}
\label{lemma-pullback-module}
Let $A \to B$ be a flat local homomorphism of local rings.
Then for any $A$-module $M$ we have
$$
\text{length}_A(M) \text{length}_B(B/\mathfrak m_AB)
=
\text{length}_B(M \otimes_A B).
$$
Here products of lengths use $0\cdot\infty=\infty\cdot0=0$.
In particular, if $\text{length}_B(B/\mathfrak m_AB) < \infty$
then $M$ has finite length if and only if $M \otimes_A B$ has finite length.
\end{lemma}

\begin{proof}
Retain the original flat local homomorphism $\alpha:A\to B$.
It is faithfully flat by Lemma \ref{lemma-local-flat-ff}.
The original closed fibre
$F=B/\mathfrak m_AB$ is nonzero, because locality implies
$\mathfrak m_AB\subset\mathfrak m_B\neq B$.
Thus $d=\operatorname{length}_B(F)$ lies in
$\mathbf Z_{\geq1}\cup\{\infty\}$.
When $M=0$, also $M\otimes_A B=0$, and the formula reads
$0=0$ using $0\cdot\infty=0$ if needed.

For every strict chain $0=M_0\subsetneq\cdots\subsetneq M_n=M$,
flatness gives injections of the actual modules $M_i\otimes_A B$
into $M\otimes_A B$ and exact sequences
$$
0\longrightarrow M_{i-1}\otimes_A B
\longrightarrow M_i\otimes_A B
\longrightarrow (M_i/M_{i-1})\otimes_A B
\longrightarrow0.
$$
Every quotient on the right is nonzero by faithful flatness,
so every step remains strict. Therefore infinite length of $M$
implies infinite length after base change, as required by
$\infty\cdot d=\infty$ for the positive value $d$.

If $0<\ell=\operatorname{length}_A(M)<\infty$, choose a
composition chain of length $\ell$. Since $A$ is local, each
simple factor is isomorphic to the original $A/\mathfrak m_A$.
The exact sequences above identify each new factor with $F$.
In the tensor order being used, define
$\Phi:(A/\mathfrak m_A)\otimes_A B\to F$ and
$\Psi:F\to(A/\mathfrak m_A)\otimes_A B$ by
$$
\begin{aligned}
\Phi((a+\mathfrak m_A)\otimes b)&=\alpha(a)b+\mathfrak m_AB,\\
\Psi(b+\mathfrak m_AB)&=(1+\mathfrak m_A)\otimes b.
\end{aligned}
$$
The first map respects tensor balancing and annihilates the
indicated representatives from $\mathfrak m_A$.
For the second map, every summand $\alpha(a)b$ with
$a\in\mathfrak m_A$ maps to
$(a+\mathfrak m_A)\otimes b=0$. Thus both maps are well
defined, and tensor balancing verifies that their composites
are identities on each pure tensor and quotient class.
Additivity along the filtration obtained by tensoring gives
$\operatorname{length}_B(M\otimes_A B)=\ell d$,
also when $d=\infty$, since a positive finite number of
infinite summands has infinite sum.
Together with the zero and infinite cases this proves the formula.

In particular, without any advance finiteness assumption on $d$,
the exact finite-length criterion is
$$
\begin{aligned}
&\operatorname{length}_B(M\otimes_A B)<\infty\\
&\quad\Longleftrightarrow\quad
M=0\ \text{or}\quad
\bigl(\operatorname{length}_A(M)<\infty\ \text{and}\ d<\infty\bigr).
\end{aligned}
$$
The positivity of $d$, faithful flatness and the proved equality
give both directions. When $d<\infty$, this specializes to
the stated equivalence.
\end{proof}

\begin{lemma}
\label{lemma-pullback-transitive}
Let $A \to B \to C$ be flat local homomorphisms of local rings. Then
$$
\text{length}_B(B/\mathfrak m_A B)
\text{length}_C(C/\mathfrak m_B C)
=
\text{length}_C(C/\mathfrak m_A C)
$$
\end{lemma}

\begin{proof}
Write the original maps as $\alpha:A\to B$ and $\beta:B\to C$.
Apply Lemma \ref{lemma-pullback-module} to $B\to C$ and the
original module $B/\mathfrak m_AB$. The receiving module in
that formula is identified with $C/\mathfrak m_AC$ by the maps
$\Phi:(B/\mathfrak m_AB)\otimes_B C\to C/\mathfrak m_AC$
and $\Psi:C/\mathfrak m_AC\to(B/\mathfrak m_AB)\otimes_B C$:
$$
\begin{aligned}
\Phi((b+\mathfrak m_AB)\otimes c)&=\beta(b)c+\mathfrak m_AC,\\
\Psi(c+\mathfrak m_AC)&=(1+\mathfrak m_AB)\otimes c.
\end{aligned}
$$
Here $\mathfrak m_AC$ is generated by the actual elements
$\beta(\alpha(a))$ for $a\in\mathfrak m_A$.
The first map respects balancing and sends any member of
$\mathfrak m_AB$ to that ideal in $C$.
For the second map, a generator times $c$ maps to
$(\alpha(a)+\mathfrak m_AB)\otimes c=0$.
Thus both are well defined, and their composites are identities
by the balancing relation. The base-change length formula now
gives precisely the claimed equality. Each of the three original
closed fibres is nonzero because the maps are local. Their
lengths are therefore positive or infinite, so their product
has no zero-times-infinity ambiguity.
\end{proof}











